\section{Discusión}
\label{sec:discusion}

\subsection{Interpretación de los resultados frente a los objetivos}

Los resultados obtenidos permiten evaluar el cumplimiento de cada objetivo
específico y, por agregación, responder la pregunta de investigación.

Respecto de \textbf{OE1}, la caracterización del ruido se materializó en
matrices de covarianza con unidades explícitas y fundamento físico trazable, y
su corrección quedó verificada por la coincidencia entre el error de la
telemetría cruda y la varianza inyectada. Respecto de \textbf{OE2}, el modelo en
espacio de estados no solo se formuló sino que se sometió a una prueba
estructural: la verificación de observabilidad del par
$(\mathbf{A}, \mathbf{H})$, que constituye la condición matemática que habilita
todo lo demás. Respecto de \textbf{OE3}, el módulo se implementó como
componente desacoplado y su corrección numérica se verificó mediante la
convergencia de la ganancia y la covarianza hacia el régimen permanente.
Respecto de \textbf{OE4}, la validación arrojó una reducción del MSE del
\SI{89,57}{\percent}, estadísticamente significativa y robusta al supuesto
distribucional.

La respuesta a la pregunta de investigación es, por tanto, afirmativa en el
alcance declarado: el diseño y la validación mediante simulación estocástica de
un módulo basado en el Filtro de Kalman Lineal permiten mitigar el ruido
instrumental y reconstruir de forma confiable el estado dinámico interno de
volumen. Ahora bien, el significado de esa respuesta requiere ser interpretado
con precisión, y a ello se dedican los apartados siguientes.

\subsection{El hallazgo central: la variable que nadie mide}

El resultado de mayor densidad conceptual no es la reducción del error sobre el
volumen sino la reconstrucción de la tasa de llenado. La distinción importa.

Sobre el volumen, el módulo mejora una señal que \emph{ya existe}: hay un sensor
que la mide, y el filtro reduce su varianza. Es una contribución valiosa pero
incremental. Sobre la tasa de llenado, en cambio, el módulo produce información
que \emph{no existía en ninguna parte del sistema}. Ningún instrumento la mide,
y el único procedimiento disponible para un SCADA convencional —la derivada
numérica de la señal cruda— entrega un error de
\SI{3,69}{\metre\cubed\per\second} sobre una magnitud cuyo valor nominal es
\SI{0,8}{\metre\cubed\per\second}. Es decir, no entrega una estimación
imprecisa: entrega ruido con la apariencia de una medición.

Que el módulo reduzca ese error a \SI{0,067}{\metre\cubed\per\second} significa
que la capa analítica no está limpiando datos sino \emph{creando observabilidad}
donde la instrumentación no la provee. Esta es la evidencia empírica más directa
de la tesis del proyecto, y es también el argumento más fuerte a favor de
concebir el estimador como componente arquitectónico y no como filtro de
preprocesamiento: un preprocesador mejora lo que recibe, mientras que este
módulo entrega a las capas superiores una variable de estado que antes no
estaba disponible para la decisión.

\subsection{Contraste con la literatura}

Los hallazgos dialogan con los tres cuerpos de literatura revisados.

Frente a la literatura de \textbf{gemelos digitales}, los resultados dan
sustento cuantitativo a la distinción conceptual de \textcite{sepasgozar2021}
entre sombra digital y gemelo digital. Si la telemetría cruda de la tasa de
llenado es esencialmente ruido, una réplica virtual alimentada directamente por
el SCADA no está representando el activo sino su distorsión. La diferencia entre
ambos conceptos deja de ser terminológica y adquiere una magnitud medible.

Frente a la literatura de \textbf{gemelos digitales mineros}, el estudio se
alinea con el diagnóstico de \textcite{ebad2026} sobre la escasa validación
industrial del área y sobre la persistencia de los problemas de integración de
datos, y con la constatación de \textcite{qu2023} sobre la adopción fragmentada
en la industria de minerales. El aporte del proyecto consiste en atender
precisamente la capa que ambas revisiones señalan como desatendida: no la
sofisticación del modelo predictivo, sino la garantía de consistencia de sus
entradas.

Frente a la literatura de \textbf{estimación de estados en procesamiento de
minerales}, los resultados son consistentes con lo reportado por
\textcite{apelt2002} en molienda semiautógena y por \textcite{barzaga2025} en
lixiviación de cobre: variables internas no instrumentadas de procesos mineros
son efectivamente reconstruibles mediante estimación de estados. La diferencia
del presente trabajo respecto de ambos no es de resultado sino de encuadre.
Aquellos estudios formulan la estimación como instrumento de medición
inferencial o de control de procesos; este la formula como componente de una
arquitectura de gemelo digital, con las exigencias de desacoplamiento,
reproducibilidad y validación estadística que ese encuadre impone.

Cabe señalar que \textcite{barzaga2025} emplean variantes no lineales —filtros
extendido y \emph{unscented}— mientras que este proyecto se restringe
deliberadamente al caso lineal. Esa restricción no es una limitación
metodológica encubierta sino una decisión de alcance: el balance de masa de una
tolva es genuinamente lineal en el rango operacional considerado, y emplear un
estimador no lineal sobre un sistema lineal añadiría costo computacional sin
ganancia de precisión.

\subsection{El hallazgo sobre robustez y su implicancia práctica}

El análisis de sensibilidad produjo un resultado que no estaba previsto en el
plan original y que tiene consecuencias prácticas relevantes. La estructura
diagonal de la malla de desajuste muestra que el desempeño del estimador está
gobernado por la \emph{razón} entre las covarianzas y no por sus valores
absolutos: un error de calibración de dos órdenes de magnitud en ambas matrices
simultáneamente degrada el RMSE desde \num{0,72} hasta apenas
\SI{0,85}{\metre\cubed}, mientras que el mismo error aplicado a una sola de
ellas lo eleva hasta \SI{8,45}{\metre\cubed}.

La implicancia para una eventual implementación industrial es directa. La
dificultad práctica de sintonizar un filtro de Kalman en planta no reside en
determinar los valores absolutos de $\mathbf{Q}$ y $\mathbf{R}$ —tarea que
requeriría caracterización metrológica detallada— sino en estimar correctamente
su proporción, es decir, cuánto confiar en el modelo físico respecto del sensor.
Esta es una tarea considerablemente más accesible para un ingeniero de procesos,
y reduce la barrera de adopción del módulo.

El resultado del barrido de degradación apunta en la misma dirección: el aporte
del módulo crece con la hostilidad del entorno, desde el
\SI{83,32}{\percent} en condiciones benignas hasta el \SI{96,68}{\percent} en
degradación extrema. Cuanto peor es el entorno minero, más necesario resulta el
componente, lo que refuerza la pertinencia contextual del proyecto.

\subsection{Limitaciones del estudio}

La interpretación de los resultados exige declarar con precisión sus límites.

La limitación principal es que la validación se realiza sobre \textbf{datos
sintéticos generados por el mismo modelo que el filtro utiliza}. Bajo esa
condición, el Filtro de Kalman es el estimador óptimo por construcción teórica,
de modo que el estudio verifica la correcta implementación y cuantifica la
magnitud de la mejora, pero no somete al estimador a una discrepancia de modelo
genuina. El análisis de desajuste de covarianzas mitiga parcialmente esta
objeción, pero no la elimina: un proceso real puede presentar no linealidades,
ruido correlacionado o de color, y derivas no modeladas que este diseño no
reproduce.

La segunda limitación es de \textbf{alcance del modelo físico}. El balance de
masa implementado es una simplificación: no incorpora la geometría real de la
cámara, la segregación granulométrica del material, ni la dependencia de la tasa
de extracción respecto de la dureza instantánea del mineral. Se excluyeron
además los estados transitorios de arranque y parada y los eventos de
atascamiento, que son precisamente las condiciones de mayor interés operacional.

La tercera limitación es la \textbf{ausencia de validación con telemetría
real}, consecuencia directa de las restricciones de confidencialidad
documentadas. Esto sitúa el trabajo en un nivel de madurez tecnológica de
laboratorio y simulación, y ninguna de las cifras reportadas puede interpretarse
como una predicción del desempeño en planta.

Una cuarta consideración, de orden estadístico, es que el rechazo del supuesto
de normalidad por la prueba de Shapiro-Wilk, si bien queda cubierto por el
contraste de Wilcoxon, sugiere que la distribución de las diferencias presenta
una estructura que este estudio no exploró en profundidad y que podría ser
objeto de análisis posterior.

\subsection{Implicancias y proyecciones}

Pese a las limitaciones declaradas, el estudio sostiene una afirmación que
excede su caso particular: en arquitecturas de inteligencia artificial
industrial, la capa de consistencia física de los datos no es una etapa de
preprocesamiento sino un componente de diseño de primer orden, y su ausencia
compromete la validez de todo lo que se construya encima.

La evidencia de extensibilidad del núcleo del módulo reportada en la
sección~\ref{sec:extensibilidad}, donde el mismo estimador opera sin
modificación sobre potencia, temperatura y desgaste, sugiere que esta capa puede
concebirse como un servicio transversal del gemelo digital antes que como una
solución puntual por variable, línea que se recoge en las proyecciones del
estudio.
